Нормализованная документная модель MongoDB хранит те же данные, но разделяет
сущности по коллекциям и связывает их ссылками \texttt{\$ref}/\texttt{\$id}.
Для этой модели используются коллекции:

\begin{itemize}
  \item \texttt{areas} --- территории;
  \item \texttt{zones} --- зоны со ссылкой на территорию;
  \item \texttt{cameras} --- камеры со ссылкой на зону;
  \item \texttt{events} --- события со ссылкой на зону;
  \item \texttt{event\_cameras} --- связи событий и камер;
  \item \texttt{camera\_telemetry} --- телеметрия со ссылкой на камеру.
\end{itemize}

Связь многие-ко-многим между событиями и камерами вынесена в
\texttt{event\_cameras}. Это уменьшает дублирование сведений о камерах по
сравнению с вложенной моделью.

Составные поля при этом остаются документами: \texttt{camera.position},
\texttt{camera.settings}, \texttt{event.payload} и
\texttt{camera\_telemetry.metrics}. Поэтому модель сопоставима с PostgreSQL
JSONB: связи нормализованы, а вариативные параметры события и телеметрии
остаются вложенными структурами.

Для аналитических запросов такая схема требует дополнительных переходов по
ссылкам. Например, чтобы сгруппировать события по территории и зоне, запросу
нужно восстановить контекст через \texttt{zones} и \texttt{areas}. Это и
создает основную стоимость нормализованной MongoDB-модели в эксперименте.
